\section{Discussion}
\label{sec:discussion}

\subsection{Relation to prior work}

The estimation problem this paper addresses, recovering disease
prevalence and code error rates without a perfect reference, has a
substantial prior literature, and we want to be precise about what is
and is not new.

\paragraph{Prevalence correction and latent-class diagnostic-test
models.} The deconvolution \eqref{eq:deconvolve} that supplies the
chart-review-calibrated prevalence estimator is the
\citet{rogan1978estimating} formula, a textbook screening-test
correction: given known sensitivity and specificity, apparent
prevalence is inverted to true prevalence. Rogan and Gladen presume
the two error rates known precisely because, as \cref{thm:glass-ceiling}
makes explicit, the triple $(\pi, \mathrm{sens}, \mathrm{spec})$ is
non-identifiable from code data alone. \citet{hui1980estimating} is
the foundational latent-class work: it estimates the sensitivity and
specificity of diagnostic tests, and the prevalence of disease, by
treating the true disease status as a latent class and fitting
multiple imperfect tests jointly. The estimator of
\cref{thm:identifiability-chart} is, in the case of two coarse features
and no chart review, a Hui-Walter latent-class model. We do not claim
to originate latent-class estimation of disease status, nor the
estimation of sensitivity and specificity without a gold standard, nor
the Rogan-Gladen plug-in. \citet{dawid1979maximum} is the
methodological ancestor one step further back: the EM algorithm for
observer (rater) error rates, which is the same latent-variable
maximum-likelihood computation applied to disagreeing raters rather
than disagreeing tests. The coding mechanism of \cref{sec:translation}
is a rater-error model in which the rater is the EHR coding process.

\paragraph{Verification bias.} The case-mix-gap residual bound
(\cref{cor:casemix}) belongs substantively to the verification-bias
literature in diagnostic-test evaluation. \citet{begg1983assessment}
showed that when verification of disease status is selective (so the
chart-reviewed or biopsy-confirmed subsample is not a random sample of
those screened), apparent sensitivity and specificity are biased, and
the bias propagates into prevalence and predictive-value estimates;
this is the now-standard verification-bias correction
\citep{hubbard2020outcome}. The substantive content of
\cref{cor:casemix} is a verification-bias statement: the subsample
sensitivity is not transportable to the cohort when chart review is
selective. What the coarsening framing contributes is positioning:
verification bias becomes a C2 failure on the chart-reviewed subsample
(non-informative coding fails within the reviewed group as it does
within the cohort), and the singleton candidate sets supplied by chart
review identify the cohort coding model only insofar as the reviewed
case mix matches. The contribution of this paper is the recasting and
the unification, not the verification-bias result itself.

\paragraph{Electronic phenotyping efforts.} eMERGE
\citep{gottesman2013emerge,newton2013validation} and the PheKB catalog
\citep{kirby2016phekb} are large phenotyping programs that build and
share code-based phenotype algorithms. PheWAS \citep{denny2010phewas}
treats code groupings as phenotypes directly.
\citet{banda2018phenotyping} surveys the methodological landscape from
rule-based to machine-learning phenotyping; \citet{liao2015highthroughput}
develop high-throughput algorithms with NLP; PheNorm
\citep{yu2018phenorm} and PheCAP \citep{zhang2019phecap} are
semi-supervised methods that reduce the labeling burden;
\citet{pivovarov2015learning} fits probabilistic phenotypes
end-to-end. The anchor-and-learn framework \citep{halpern2016anchor}
learns phenotype classifiers using high-precision ``anchor''
observations whose presence almost certainly implies the true state;
an anchor is structurally a near-singleton candidate set, and
\citet{halpern2016anchor} is the closest prior work in spirit to the
singleton-restores-identifiability mechanism of this paper.
Surveillance bias, the phenomenon that more intense looking finds more
disease, is named by \citet{haut2011surveillance} and quantified as a
phenotyping artifact by \citet{carrell2014surveillance}. The
EHR-as-data perspective of \citet{hripcsak2013correlating} and the
OHDSI treatment-pathway work of \citet{hripcsak2016characterizing}
frame the data-quality and process-dependence issues that the
coarsening lens makes precise.

\paragraph{What is new here.} The contribution is not latent-class
phenotyping. It is the masked-cause unification: (i) the explicit
candidate-set translation that places code-based phenotyping inside the
coarsening-at-random framework of \citet{towell2026masked}; (ii) the
C1-C2-C3 classification that identifies informative coding as
specifically a C2 violation, distinguishing it from miscoding (C1) and
from confounding between coding parameters and prevalence (C3); (iii)
the explicit glass-ceiling construction (\cref{thm:glass-ceiling}),
which exhibits the non-identifiable solution surface rather than
asserting non-identifiability; and (iv) the informative-coding bias
bound (\cref{thm:bias-informative}) and the case-mix-gap residual bound
(\cref{cor:casemix}), which quantify the bias of the code-only and
calibrated estimators in terms of interpretable quantities, the
severity-coding correlation and the case-mix gap. The Hui-Walter model
tells a practitioner \emph{that} prevalence can be estimated from
imperfect tests; the coarsening framework tells the practitioner
\emph{which assumption} (C2) the EHR violates, \emph{why} (informative
coding), and \emph{how large} the resulting bias is.

\subsection{What the framework does not address}

\begin{itemize}
\item \textbf{Multi-code phenotypes and temporal structure.} We treat
  the candidate set as generated by a small set of binary codes
  observed cross-sectionally. Real phenotype algorithms use code counts,
  code timing, medication sequences, and laboratory trajectories. The
  framework extends to richer candidate sets, but the present analysis
  does not.

\item \textbf{Chart review error.} We treat the chart-reviewed state as
  a true singleton. Adjudication itself is imperfect, and inter-rater
  disagreement is exactly the setting of \citet{dawid1979maximum}. The
  framework would benefit from incorporating chart-review error via the
  C2-relaxation results of \citet{towell2026mdrelax}.

\item \textbf{The form of the coding mechanism.} The bias bound
  \eqref{eq:meth-bound} assumes a logistic coding mechanism with a
  scalar severity. Real coding depends on insurance, site, clinician,
  and calendar time. A more flexible misspecification model is needed.

\item \textbf{Causal estimands.} We target descriptive cohort
  prevalence. Phenotyping feeds downstream association and
  causal-inference analyses, where informative coding induces further
  bias (differential misclassification of the exposure or outcome). The
  propagation of the coarsening bias into those estimands is left open.

\item \textbf{Identifiability from codes alone via parametric
  restriction.} \Cref{thm:glass-ceiling} concerns an unrestricted
  coding mechanism. With a sufficiently restricted parametric family,
  or multiple conditionally independent codes as in
  \citet{hui1980estimating}, prevalence can be identified without chart
  review. Characterizing exactly which restrictions suffice, and how
  fragile the resulting estimates are, is a natural follow-up.
\end{itemize}

\subsection{Sibling applications of the framework}

The masked-data lens has now been applied across a small set of
measurement domains.
In single-cell RNA sequencing, \citet{towell2026scrnacoarsening} treats
dropout zeros as candidate-set ambiguity and identifies ERCC
spike-ins as the singleton mechanism. In spatial transcriptomics,
\citet{towell2026spatialcoarsening} ports the rank condition to
cell-type deconvolution within spots. In differential privacy,
\citet{towell2026dpcoarsening} treats noisy release as the masking
mechanism and characterizes identifiability of population parameters
under privacy budgets. In programmatic weak supervision,
\citet{towell2026weaksupcoarsening} treats labeling functions as
candidate sets and develops identifiability conditions in the absence
of a gold standard. Electronic phenotyping joins this family as the
healthcare-data instance: the same masked-cause structure with the
same singleton-restores-identifiability mechanism, distinguished by
informative coding as the dominant C2 violation rather than (as in
the transcriptomics applications) a comparatively mild deviation from
non-informative observation. Recognition of this common structure may
streamline methodology across these subfields.
